%% notes-tex/eqtl-data-model/sections/2-observed-vs-inferred.tex
%% [[torchcell.datamodels.eqtl-data-model]]
%% https://github.com/Mjvolk3/torchcell/tree/main/notes-tex/eqtl-data-model/sections/2-observed-vs-inferred.tex
%% SOURCE: hand-authored; quotes are verbatim from the mirrored
%% $DATA_ROOT/torchcell-library/boocockSinglecellEQTLMapping2025/paper.md (lines 50, 71, 208, 212, 228).
%% The error-budget arithmetic is derived from published parameters and is labeled as such.

\section{The genotype is inferred, not observed\secstatus{todo}}\label{sec:observed}

\subsection{What the methods actually say\secstatus{todo}}

Verbatim from the mirrored Boocock 2025 paper:

\begin{quote}
\small
``We used a hidden Markov model (HMM) to reconstruct the patterns of inheritance of
parental alleles in each cell based on the observed genotypes at expressed SNPs.''
\end{quote}

\begin{quote}
\small
``We defined prior genotype probabilities as 0.5 for each parental variant. Emission
probabilities were calculated as previously described for low coverage sequencing data.''
\end{quote}

\begin{quote}
\small
``Genotype probabilities were standardized, and markers in very high LD ($r > 0.999$) were
pruned. This LD pruning is approximately equivalent to using markers spaced 4 centimorgans
(cm) apart.''
\end{quote}

Two terms in that last quote: \term{LD}, linkage disequilibrium, is the statistical
correlation between alleles at different sites, near-perfect for neighboring sites here
because they are inherited together in blocks; a \term{centimorgan} (cM) is a unit of
genetic map distance, two loci sitting 1 cM apart when 1\% of meioses place a crossover
between them.

So the genotype matrix is not binary. It is

$$
X \in [0,1]^{n \times m}, \qquad
X_{ij} \;=\; \Pr\!\big(\text{segregant } i \text{ carries the RM allele at } j \,\big|\, \text{reads}\big)
$$

and the eQTL model is fit on the standardized posterior directly. Binarizing at $0.5$ is
done only to compare cell identities.

\subsection{Three levels of partial observation\secstatus{todo}}

\begin{table}[H]\centering
\small
\caption[]{What is actually observed, by study. The single-cell design is the most severely
partial: genotype coverage is conditioned on the transcript being expressed, so the
genotype observation and the phenotype measurement come from the same reads.}
\label{tab:levels}
\begin{tabular}{@{}l l l@{}}
\toprule
Study & Observed & Genotype is \\
\midrule
Bloom 2013 \citep{bloomFindingSourcesMissing2013} & low-coverage WGS of each segregant & calls at $\sim$11,623 markers \\
Nguyen Ba 2022 \citep{nguyenbaBarcodedBulkMapping2022} & low-coverage WGS plus barcode & calls at markers, $n \sim 10^5$ \\
Boocock 2025 \citep{boocockSinglecellEQTLMapping2025} & only SNPs that happen to be expressed & an HMM posterior per marker \\
\bottomrule
\end{tabular}
\end{table}

The observation channel differs per study, so the probabilistic genotype has to be
extracted differently per study: there is no shared genotype-recovery path for a loader
to reuse, only a shared representation for its output.

A gene that is switched off in a cell contributes no genotype information about its own
locus. That coupling has no analog in any other dataset in the collection, it is inherent
to the one-pot design, and the record can only carry it, not correct it.

\subsection{Why sparse observation is nevertheless enough\secstatus{todo}}\label{sec:enough}

The segregant is not being sequenced de novo. It is being \emph{assigned} to one of two
known genomes at each locus, which is a far easier problem:

\begin{enumerate}[leftmargin=1.4em,itemsep=1pt,topsep=2pt]
  \item The parents are deep-sequenced, ``deep ($>100\times$) paired-end sequenced data for
    our parental strains,'' so the complete variant catalog and both alleles at every site
    are known before any segregant is looked at.
  \item Meiosis produces few boundaries. In \org{S.\ cerevisiae} roughly 90 crossovers per
    meiosis \citep{manceraHighResolutionMapping2008}, so across 16 chromosomes a segregant
    genome is on the order of $10^2$ contiguous blocks.
  \item Within a block, parental origin is constant, so one marker fixes every cataloged
    variant in its block.
\end{enumerate}

Reconstruction reduces to locating $\sim\!10^2$ boundaries rather than determining
$\sim\!10^4$ independent values.

\notesec{Error budget, derived from published parameters rather than measured} BY and RM
segregate at $\sim\!4.5\times10^4$ SNPs over 12 Mb, about 1 per 260 bp. Bloom 2013's 11,623
markers give $\sim\!1$ kb spacing, so a crossover is localized to about 4 variants:

$$
\underbrace{90}_{\text{crossovers}} \times \underbrace{4}_{\text{variants per interval}}
\;\approx\; 360 \;\text{ of } 4.5\times10^4 \;\approx\; 0.8\%\ \text{ambiguous}
$$

leaving roughly $99\%$ of cataloged variant sites assignable. \external{arithmetic from
published crossover and variant-density figures, not a measured concordance}

\notesec{The larger error source is not crossovers} Yeast meiosis also produces $\sim\!46$
non-crossover \term{gene conversion} events per meiosis with tracts of order 2 kb
\citep{manceraHighResolutionMapping2008}. A gene conversion is a byproduct of meiotic
recombination repair: a double-strand break is repaired by copying sequence from the
homologous chromosome, so a short tract of one parent's sequence is overwritten with the
other parent's, without exchanging the flanking arms. The result is a $\sim\!2$ kb patch
of the ``wrong'' parent inside an otherwise intact block. Those patches are short, they do
not shift the surrounding haplotype, and sparse markers miss them entirely, so they affect
more base pairs than boundary uncertainty and are harder to detect.

\subsection{What reconstruction cannot give you at any marker density\secstatus{todo}}

\begin{itemize}[leftmargin=1.2em,itemsep=1pt,topsep=2pt]
  \item \textbf{De novo mutations.} A segregant is a mosaic plus whatever arose since the
    cross. Few per strain, invisible to a parental-assignment method, so the genome is
    never exactly a mosaic.
  \item \textbf{Aneuploidy and copy number.} Haploidy does not preclude these: a haploid
    segregant can carry a second copy of a whole chromosome (a disomy) or a segmental
    amplification, and yeast segregants do acquire them. Detecting that needs read
    coverage, not marker identity. The mirrored Boocock 2025 text reports no aneuploidy
    or copy-number screen, so for that design it is an unmodeled error term rather than a
    measured one.
  \item \textbf{Non-reference sequence.} A variant call file is not an assembly.
    Reconstructing total genomic content needs the parents' assemblies, including sequence
    absent from S288C. The ingredients exist: the deep paired-end parental data supports
    de novo assembly, with a known residual error rate, so holding pinned parent
    assemblies is an attemptable requirement rather than a new sequencing campaign.
\end{itemize}

\subsection{Per cell the answer flips\secstatus{todo}}

Boocock reports \emph{median genotype agreement of 92.5\%} between the single-cell HMM
genotypes and whole-genome sequencing of the same strains, improving with UMI count. About
1 marker in 13 wrong is entirely acceptable for an association fit across 27,744 cells,
where the posterior enters as a standardized covariate and the error averages out. It is
nowhere near sufficient to write down a genome. What rescues the strain-level genotype is
pooling: a median of 17 cells per segregant, against 393 segregants that had already been
whole-genome sequenced in Bloom 2013 and against which the cells were matched. The record
never needs the per-cell genotype to be a genome; what it stores is the strain-level
mosaic, and that is the same compositional form the schema already uses everywhere else,
a pinned reference plus recorded differences, so the insufficiency of a single cell's
genotype does not threaten the ingestion pattern.

\notesec{Do not conflate reconstruction resolution with mapping resolution} The HMM used
all expressed SNPs. The 4 cM LD-pruned marker set was used for the association, because
markers above $r>0.999$ carry no independent information. Coarse mapping resolution says
nothing about how well the genome can be reconstructed.

\subsection{Reconstruction and attribution come apart\secstatus{todo}}\label{sec:verdict}

Two questions about the same data sound alike and have opposite answers.

\notesec{Reconstruction succeeds} \emph{Which parent's allele does segregant $i$ carry at
site $j$?} This is a two-way choice between two fully sequenced genomes, decided jointly
for a whole block by every marker in it, and the arithmetic of
Sec.~\ref{sec:enough} puts it at $\sim\!99\%$ of sites. More segregants, denser markers or
deeper coverage all push it further toward certainty.

\notesec{Attribution is blocked} \emph{Which variant causes the expression difference?}
Within a haplotype block every variant is inherited with every other, so their columns in
$X$ are identical and no fit can tell them apart. The effect is identifiable only up to
the block. Collecting more segregants does not help, because the collinearity comes from
the cross itself: the same meioses that made the panel made the blocks. Boocock's pruning
of markers above $r>0.999$ concedes exactly this point; markers that co-segregate carry no
separate information.

\notesec{What each answer decides} Reconstruction is what the sequence gate on the
candidate list asks about: can the total genomic content of the strain be estimated? It
can, so \textbf{segregant panels pass}, conditional on holding the parent assemblies.
Attribution is what modeling wants and does not get: a cross localizes an effect to an
interval and nothing more. Getting from an interval to a base takes an assay that tests
alleles one at a time, such as a massively parallel reporter assay, which is why those
datasets belong in the same ingestion wave rather than a later one.
